\section{Results: controls, implementations and defenses that hold}
\label{sec:results2}

Section~\ref{sec:results} reports a negative result about one family of
defenses. This section reports three things that complicate it: controls that
changed our own conclusions, a test of whether a result belongs to a principle
or to one codebase, and the defenses that did transfer.

\subsection{Three controls that changed our own conclusions}

Each of the following would have produced a wrong published claim if we had
reported the metric the relevant literature reports.

\textbf{A no-retrieval baseline.} DyPRAG~\cite{tan-2025-dynamic} is perfectly
flat under an attack that destroys the undefended arm, scoring 5.8, 6.0, 5.2
and 6.0 across the dose range while the undefended arm falls from 63.8\% to
2.5\%. Read alone, that is the strongest defense result in the campaign. The
closed-book arm scores 7.5\% and is statistically indistinguishable from it.
\textbf{DyPRAG is flat because it is not reading the document}, not because it
is resisting the attack. Only a no-retrieval arm separates immunity from
inattention, and no paper in this family reports one.

\textbf{An abstention axis.} COIECD is inert on accuracy at every dose. Dumping
and classifying every generation shows it is not inert at all. On the items
where the undefended model answered from memory, COIECD converts
\textbf{32.1\%} of those answers into explicit abstentions, and does so only
under attack; on clean context the same pairing moves almost nothing. Because
the parametric answer is wrong by construction in this data, what it suppressed
were confident wrong answers and what it produced instead was "not mentioned".
That is a real safety gain, and it is invisible in the accuracy column that
every paper in this family reports.

\textbf{A model control.} RobustRAG ran on Mistral and the context-reliance
family ran on Llama, because those are the models each paper released. Re-running
RobustRAG on the other family's model changes outcomes. The tone-steering
finding survives unchanged, which strengthens it. A second finding inverts.
\textbf{Holding defense and attack fixed and changing only the base model
changes an effect's size, its existence, and in one case its sign.}

\subsection{Principle or implementation}

A single implementation cannot tell you whether you have found a property of a
design principle or a quirk of one codebase. RobustRAG~\cite{xiang-2024-certifiably}
ships two aggregation variants that share the isolate-then-aggregate principle
and share no decision logic downstream of it. One extracts keywords and
re-prompts; the other averages next-token distributions.

Both show the same behaviour under tone steering, at or below 0.3\% across four
cells and two models, against CK-PLUG's 99.0\% on the same attack and model.
\textbf{Aggregation does not open the tone-steering hole. Context-reliance
does.} That the two independent implementations agree is what licenses stating
it as a property of the principle.

The caution runs the other way too. Across fifteen dose levels the two
implementations agree within noise on thirteen. On one they do not. At nine
poisoned passages in ten, against refusal induction, KeywordAgg is
\textbf{+12.3\pp} and DecodingAgg is \textbf{-17.3\pp}, a gap of 29.7 points
with both ends clearing the 11.4\pp threshold at $n=300$. One implementation of
isolate-then-aggregate is resolvably protective and the other is resolvably
harmful, on the same attack, at the same dose, inside one codebase.

The gap grows as the honest minority of passages shrinks, from 13.0 points at
five in ten to 29.7 at nine in ten, and then collapses to 1.7 at ten in ten
when no honest passage remains and the two implementations have nothing left to
disagree about. An SoK that reported one number per published method would have
missed all of this.

\subsection{Defenses that do transfer, and one line that fixes one that does not}

The picture is not uniformly negative, and the exceptions are informative.

RobustRAG stops PoisonedRAG, which its own paper tested, and we reproduce that
(+11.0\pp). It then fails completely against corpus poisoning at displace-10,
scoring 4.6\% against the undefended arm's 3.4\%. The reason is not its
security mechanism. Corpus poisoning displaces every honest passage out of the
top-$k$ window, so isolate-then-aggregate has nothing honest left to aggregate.
Backfilling the window to its intended size restores \textbf{45.2\%}, a gain of
40.6\pp. \textbf{The defense was sound and the retrieval interface was not.}

Across the agent cell, 24 of 56 conditions show the technique holding with
attack success driven to zero, 4 show a resolvable reduction, 19 show no
resolvable effect at that sample size, and 9 show no reduction at all.
Ingestion-stage techniques account for most of the holds.

\subsection{The dose where a defense has its best chance}
\label{sec:lowdose}

Results so far are reported at saturation, where the planted passage replaces
the entire context. That corner is where the attack is strongest, and it is
also where the measurement is weakest: with the real document entirely gone, no
method can read it, so a low accuracy partly measures the attack being complete
rather than the defense failing. Carlini makes the equivalent objection to
DeepSec's distortion budgets~\cite{carlini-2019-critique}, and it applies to us.

So we report the other end as well. At the \textbf{10\% dose} one chunk in ten
is replaced and nine tenths of the honest document remain. If a method can
discount a planted passage while still using a good one, this is where it
shows.

\begin{table}[t]
\centering
\caption{Every defense against every payload it met, at the 10\% dose. ``Helps''
means the defended arm beats its own undefended arm by more than that run's
detectability threshold. Worst case is the payload with the largest shortfall.}
\label{tab:lowdose}
\small
\begin{tabular}{lrrl}
\toprule
defense & payloads & helps on & worst payload met \\
\midrule
ACD       & 4 & 0 & corpus poisoning \;\; $-3.2$ \\
CK-PLUG   & 3 & 0 & corpus poisoning \;\; $-14.9$ \\
COIECD    & 3 & 0 & distractor \;\; $-0.8$ \\
DyPRAG    & 1 & 0 & corpus poisoning \;\; $-30.2$ \\
ParamMute & 3 & 0 & distractor \;\; $-1.0$ \\
SCR       & 1 & 0 & corpus poisoning \;\; $-12.8$ \\
\midrule
\textbf{total} & \textbf{15} & \textbf{0} & \\
\bottomrule
\end{tabular}
\end{table}

\textbf{Not one of the 15 (defense, payload) pairs shows a resolvable benefit
at the 10\% dose.} Every defense is at or below its own undefended arm on every
payload it met, at the dose most favourable to it. The saturation numbers
elsewhere in this paper are the dramatic version of the result;
Table~\ref{tab:lowdose} is the defensible one.

The table is also the worst-case reading that a security claim requires. A
defense beaten by one of three payloads is not two-thirds robust, it is broken,
and averaging across payloads would hide that. We report the minimum rather
than the mean for the same reason Carlini gives: robustness is a property of
the strongest attack a method has met, not of the attacks its authors chose.

\subsection{An adaptive attacker}
\label{sec:adaptive}

Every result above comes from an attack that was not built for the defense it
met. Section~\ref{sec:experiments} states why that is conservative for a
failure and invalid for a resistance. ACD is our one resistance result, so it
is the one that has to be tested properly.

ACD mixes the no-context and with-context next-token distributions by
\[
d_\alpha = \frac{H}{H_{\mathrm{ctx}} + H},
\qquad
s \leftarrow s + d_\alpha\,(s_{\mathrm{ctx}} - s),
\]
where $H$ is the Shannon entropy of the no-context distribution and
$H_{\mathrm{ctx}}$ that of the with-context one. $H$ is a property of the
question. $H_{\mathrm{ctx}}$ is conditioned on the retrieved context, which an
attacker who owns the corpus controls outright. \textbf{So the attack has
exactly one objective: minimize $H_{\mathrm{ctx}}$.}

\textbf{The ceiling follows from the rule, not from an experiment.} As
$H_{\mathrm{ctx}} \to 0$, $d_\alpha \to 1$ and the update collapses to
$s_{\mathrm{ctx}}$ unchanged, which is the undefended model. A perfect adaptive
attack of this kind turns ACD into plain RAG and can do no better; it cannot
drive ACD \emph{below} its own undefended arm, because $d_\alpha$ is bounded in
$(0,1)$ and the endpoints are closed-book and plain RAG. We committed that
bound, and five further predictions, before building the payload. An ACD arm
landing below the undefended arm would have falsified our account of the
mechanism.

The attack is white-box and query-based: five candidate passages per item, all
asserting the \emph{same} distractor answer and differing only in form, each
scored by the $H_{\mathrm{ctx}}$ ACD itself would compute, keeping the
minimum. Holding the asserted answer fixed is what makes any movement
attributable to form rather than to a different claim. It is a deliberately
weak optimizer, so what it achieves is again a lower bound.

\begin{table}[t]
\centering
\caption{ACD against the distractor payload in its plain form and in the form
that minimizes ACD's own entropy term, at saturation, $n=686$. The adaptive
payload is stronger against the undefended arm too, so the defensive gap is
the quantity that isolates the defense.}
\label{tab:adaptive}
\small
\begin{tabular}{lrrr}
\toprule
at saturation & plain & adaptive & change \\
\midrule
median $H_{\mathrm{ctx}}$ & 1.348 & 0.783 & $-42\%$ \\
median $d_\alpha$ & 0.525 & 0.653 & $+0.13$ \\
\midrule
undefended copies attacker & 76.0 & 89.1 & \\
ACD copies attacker & 40.8 & \textbf{64.0} & $+23.2$ \\
\textbf{defensive gap} & \textbf{35.1} & \textbf{25.1} & $\mathbf{-29\%}$ \\
\midrule
ACD recalls from memory & 32.5 & 22.2 & \\
ACD answers correctly & 5.0 & 3.2 & \\
\bottomrule
\end{tabular}
\end{table}

All six predictions held. $H_{\mathrm{ctx}}$ falls 42\%, $d_\alpha$ rises
toward 1, ACD's copy rate rises from 40.8\% to \textbf{64.0\%}
($p = 9.5\times10^{-18}$), its parametric recall falls toward the undefended
arm's, and its accuracy lands on \textbf{3.2\%} --- exactly the undefended
arm, and below the no-retrieval floor of 8.2\%. ACD does not fall below the
undefended arm on copy rate, as the ceiling requires.

\textbf{The honest size of the effect is the gap, not ACD's rise.} The adaptive
payload is stronger against the undefended arm as well, so reporting only
ACD's change would credit the defense with damage the attack does to everyone.
Against its own undefended arm the protection falls from 35.1\pp{} to
25.1\pp{}: \textbf{29\% of it removed by a best-of-five search over sentence
templates.} Not all of it. $d_\alpha$ reached 0.653 of a possible 1.0, so a
gradient search on $H_{\mathrm{ctx}}$ has the remainder to work with.

What this licenses is narrow and worth stating as such. It does not show ACD is
broken. It shows that the one defense in our campaign that looked like a
counterexample degrades under an attacker who reads the method and spends five
forward passes per item, and that every published evaluation of a
context-reliance defense in our corpus is non-adaptive. The gap between
non-adaptive and adaptive evaluation is large enough to change the verdict.
